% ==================== 英文摘要 ====================
\clearpage
\phantomsection
\addcontentsline{toc}{chapter}{Abstract}
\thispagestyle{fancy}

\begin{center}
{\zihao{3}\bfseries \thesisTitleEn}
\end{center}

\vspace{0.5em}

{\zihao{-4}\setstretch{1.25}

\noindent\hspace{2em}{\bfseries Abstract: }In dense multi-obstacle scenes, mainstream autonomous driving path planners are prone to blockage and over-conservative nudging. This thesis takes the Planning module of Baidu's Apollo 11.0 open-source autonomous driving platform as the study object and dissects the path-boundary construction logic carried by PathBoundsDeciderUtil. Three deficiencies are identified. The uniform safety buffer squeezes the passage band beyond what local obstacle geometry requires. Per-obstacle nudge decisions are taken independently and may assign mutually contradictory avoidance directions. Dynamic obstacles never enter path-boundary computation, so the downstream path and speed planners cannot act consistently.

MIKU, the Multi-obstacle Interaction-density Kinematic-prediction Unified Planner, is designed against these three deficiencies and casts multi-obstacle nudge direction selection as a combinatorial optimization problem. For safety clearance, a multi-factor threat score combining time-to-collision (TTC), lateral overlap ratio, relative velocity, obstacle type, and interaction density replaces the uniform buffer, yielding a differentiated margin $\delta_i$ per obstacle. For decision granularity, a sweep-line procedure groups longitudinally overlapping obstacles into connected components and decides nudge directions jointly within each group. We prove that the optimal partition stays contiguous under lateral ordering and that the maximum-gap strategy attains the optimum, which reduces the combinatorial search from $O(2^k)$ to $O(k\log k)$. For spatiotemporal coordination, an arrival-time function $\tau(s)$ together with predicted trajectories admits dynamic obstacles into path-boundary construction, builds a spatiotemporal traversable corridor, and fuses its boundary with the Station-Time (ST) constraints emitted by SpeedBoundsDecider inside the speed Quadratic Programming (QP) layer.

Eight simulation scenarios are evaluated, \MPressureScenarioTotal{} original stress cases and \MComparableScenarioTotal{} lightly modified comparable cases. On the stress cases, Baseline passes \MPressureBaselinePassCount{}/\MPressureScenarioTotal{} and MIKU passes \MPressureMikuPassCount{}/\MPressureScenarioTotal{}. On the comparable cases both methods reach the goal in every run, while MIKU lowers longitudinal acceleration and jerk peaks in the dynamic scenarios and reduces the average total QP solving time from \MComparableBaselineQpTotal\,ms to \MComparableMikuQpTotal\,ms. The overall complexity stays at $O(n\log n+K)$, with $K$ the path or time discretization scale, which fits the real-time budget of the Planning module. On the engineering side, MIKU enters Apollo as a new path-boundary branch of LaneFollowPath and feeds its corridor constraints to PiecewiseJerkSpeedOptimizer through STDrivableBoundary.

\vspace{0.5em}
\noindent\hspace{2em}{\bfseries Keywords: }Autonomous Driving; Path Planning; Apollo; MIKU Algorithm; Combinatorial Optimization

}
